\section{A quantitative rank and branch reduction for Heisenberg factors}
\label{sec:heisenberg-rank-concavity}

\paragraph{Status and provenance.}
D-TD57, question Q2fm, refines the exact TD08--10 inequalities
using the actual original marks omitted by a projective factor.
The parent requested the strict-gap investigation and independently
reconstructed the proposed concavity minimum and endpoint constants.
The complete argument below is by \texttt{/root/imprimitive}.
The complete author proof, source-interface, code and successful-output
readings pass.  Complete root and independent whole-proof reviews also
pass; the retained formatting-check failure is documented separately.
It is a strict reduction of the remaining branches, not an unrestricted
bound on the Heisenberg exponent.

\paragraph{Actual source and primitive factor.}
Let \(C\to\mathbf P^1\) be a connected finite \(J\)-Galois cover, with
faithful irreducible coefficient representation \(V\), \(e=\dim V\).
Its actual generating product-one inertia tuple satisfies
\[
 T_i\ne I,\quad r\ge12,\quad
 \sum_{i=1}^r\operatorname{rank}(T_i-I)=2e+6,\quad
 \frac{\operatorname{rank}(T_i-I)}e\ge\frac1M
 \quad(1\le i\le r),\qquad M\ge1.
 \tag{HRC.1}
\]
Suppose \(V=W\otimes U\) is preserved factorwise, with its actual finite
projective factor image \(P\subset\operatorname{PGL}(W)\).
Assume that \(P\) is primitive and contains a normal irreducibly
represented projective Heisenberg subgroup
\[
 Q\simeq\mathbf F_p^{\,2m},\qquad m\ge1,\qquad
 a=\dim W=p^m,\qquad e=a\dim U.
 \tag{HRC.2}
\]
Its scalar commutator pairing is nondegenerate, as in TD05.
Primitivity of the original \(J\) is a sufficient condition for
the assumed primitivity of this actual factor image, by TD11.
The whole-original floor in (HRC.1) is explicit; it is not a
statement about only the selected factor.

Let \(y_i\in P\) be the projected inertia, and \(g_i\in\operatorname{Sp}(Q)\)
its actual conjugation action on the label space.  Write
\[
 I=\{i:y_i\ne1\},\qquad k=|I|,\qquad
 s_i=\operatorname{rank}(g_i-I_Q)\quad(i\in I),\qquad
 \delta_p=\frac{p-1}{2p}.
 \tag{HRC.3}
\]
The whole tuple \(g_i\) generates its actual image and has product one.
Projected identity entries have \(s_i=0\), and may be omitted from
the rank sum, but remain genuine original marks in (HRC.1).

\begin{theorem}[D-TD57]
Under these hypotheses \(k\ge3\).  If \(m\ge k-1\), then the
projective supports of any factor representatives \(A_i\) satisfy
\[
 \sum_{i\in I}\frac{\nu_W(A_i)}a
 \ge
 1+(k-2)\delta_p-\frac{1+p^{\,k-2}}{2a},
 \tag{HRC.4}
\]
where \(\nu_W(A)=a-\max_\lambda\dim\ker(A-\lambda I)\).
The actual whole-source budget consequently gives
\[
 (k-2)\delta_p-1+\frac{r-k}{M}
                  \le\frac{13+p^{\,k-2}}{2a}.
 \tag{HRC.5}
\]
The following alternatives hold without the auxiliary assumption
\(m\ge k-1\):
\[
 \begin{array}{c|c}
 p&\text{necessary alternative}\\ \hline
 2&
 a\le\max\{64,29M/12\}\quad\text{or}\quad 3\le k\le5,\\[2pt]
 3&
 a\le\max\{81,20M/7\}\quad\text{or}\quad 3\le k\le4,\\[2pt]
 p\ge5&
 m\le3\quad\text{or}\quad 3\le k\le4.
 \end{array}
 \tag{HRC.6}
\]
More precisely, at \(p=2,m\ge7\), one has \(k\le6\), and
\(k=6\) forces \(a\le29M/[2(r-6)]\).
At \(p=3,m\ge5\), one has \(k\le5\), and
\(k=5\) forces \(a\le20M/(r-5)\).
\end{theorem}

\begin{proof}
\emph{Step 1: the actual group and rank interfaces.}
The projected tuple generates \(P\) and has product one.
If it had at most two nonidentity entries, \(P\) would be cyclic:
with two, the entries are inverse to one another.  This is impossible
because \(P\) contains the noncyclic subgroup
\(Q\simeq(C_p)^{2m}\).  Thus \(k\ge3\); no lower bound of twelve
on \(k\) is used.

For every nonidentity \(y_i\), TD08 gives the all-target estimate
\(\nu_W(A_i)/a\ge\delta_p\).  The exact trace formula also gives
\[
 \frac{\nu_W(A_i)}a\ge
    \frac{1-p^{-s_i/2}}2\quad(s_i>0).
 \tag{HRC.7}
\]
Both estimates are unaffected by scalars and survive tensoring with
the entire arbitrary \(U\).  In particular they bound the actual
normalized moving rank of \(T_i\).

Let \(H=\langle g_1,\ldots,g_r\rangle\).  Primitivity of \(P\)
implies \(Q^H=0\).  Otherwise a nonzero fixed label would give
a projectively invariant order-\(p\) Weyl operator, and its \(p\)
equal-dimensional eigenspaces would be permuted by \(P\).
That would be a nontrivial imprimitivity system.
Scott's inequality for the actual product-one generating tuple,
with its symplectic self-duality, therefore gives precisely the
accepted TD10 statement
\[
                      \sum_{i\in I}s_i\ge4m.
 \tag{HRC.8}
\]
The known lemma is the arbitrary-field statement in
\cite[Lemma 2.3, author version p.~7]{TD:GuralnickMaroti2011},
with original provenance \cite{TD:Scott1977}; its full source and
fixed-space interfaces were checked in TD10.
No modular absolute irreducibility, classification, or scalar
normalizer splitting is being added here.

\emph{Step 2: the relaxed rank polytope.}
Set \(d=2m\), and for \(i\in I\) put
\[
 R_i=\max\{2,s_i\},\qquad
 f(x)=\frac{1-p^{-x/2}}2\quad(2\le x\le d).
 \tag{HRC.9}
\]
When \(s_i=0\) or1, the sharp uniform estimate is
\(\delta_p=f(2)\).  When \(s_i\ge2\), use (HRC.7).
Consequently \(\nu_W(A_i)/a\ge f(R_i)\) in every case.
Also \(2\le R_i\le d\), and (HRC.8) implies
\(\sum R_i\ge2d\).

Now assume \(m\ge k-1\).  Since \(k\ge3\), this implies
\(d\ge2k-2\), so \(2k\le2d\le kd\).
The function \(f\) is increasing and concave.  Reducing coordinates
if needed shows that the minimum over the relaxed polytope
\[
          2\le x_i\le d,\qquad \sum_{i=1}^k x_i\ge2d
 \tag{HRC.10}
\]
is attained with sum exactly \(2d\).
One can then move mass from a smaller interior coordinate to a
larger interior coordinate until an endpoint is reached.  Concavity
ensures that this does not increase the sum of the \(f(x_i)\).
Repeating gives a vector with all but at most one coordinate at
the endpoints2 and \(d\).

In detail, above the all-two vector the required excess is
\(2d-2k\).  Since \(d\ge2k-2\), it is at least \(d-2\),
and since \(k\ge3\), it is less than \(2(d-2)\).
Thus, allowing a degenerate residual endpoint, the minimizing vector is
\[
                 (d,\ d-2k+4,\ 2,\ldots,2).
 \tag{HRC.11}
\]
Its residual coordinate is at least2 and at most \(d\).
This is a relaxation of the actual symplectic rank possibilities;
no assertion that the displayed vector is realized by a generating
tuple is necessary.  Evaluating \(f\) on (HRC.11) proves (HRC.4).

\emph{Step 3: retain all omitted original marks.}
Every \(i\notin I\) is still a nonidentity original coefficient
operator and retains its cost at least \(1/M\).  Scalar tensor
ambiguity is harmless because all selected-factor estimates use
the largest eigenspace, not a chosen factor eigenvalue.
Using \(e\ge a\) gives
\[
 \begin{split}
 2+\frac6a
 &\ge2+\frac6e\\
 &\ge1+(k-2)\delta_p-\frac{1+p^{\,k-2}}{2a}
                     +\frac{r-k}{M}.
 \end{split}
 \tag{HRC.12}
\]
Rearranging is exactly (HRC.5).  The support floor alone also gives
the useful elementary upper bound
\[
                      k\delta_p\le2+\frac6e\le2+\frac6a.
 \tag{HRC.13}
\]

\emph{Step 4: every low-\(m\) endpoint in (HRC.6).}
For \(p=2\), if \(m\le6\), then \(a\le64\).
Otherwise \(m\ge7\), \(a\ge128\), and (HRC.13) implies \(k\le8\).
If \(k=7\) or8, then \(m\ge k-1\), so (HRC.5) applies.
Discarding its nonnegative original-mark term gives, respectively,
\[
        k=7:\quad a\le90,\qquad
        k=8:\quad a\le77.
 \tag{HRC.14}
\]
Both contradict \(a\ge128\).  Hence \(k\le6\).
For \(k=6\), the constant term \((k-2)\delta_p-1\) is zero,
and (HRC.5) gives
\(a\le29M/[2(r-6)]\le29M/12\).
For \(k\le5\), the first remaining branch in (HRC.6) holds.

For \(p=3\), if \(m\le4\), then \(a\le81\).
Otherwise \(m\ge5\), \(a\ge243\), and (HRC.13) implies \(k\le6\).
If \(k=6\), (HRC.5) is applicable and, after discarding the
nonnegative original-mark term, gives \(a\le141\), a contradiction.
Thus \(k\le5\).  For \(k=5\), that inequality has zero constant
term and yields \(a\le20M/(r-5)\le20M/7\).

Finally let \(p\ge5\).  If \(m\le3\), the last bounded-exponent
alternative is already stated.  Otherwise \(a\ge p^4\ge625\)
and \(\delta_p\ge2/5\), so (HRC.13) implies \(k\le5\).
For \(k=5\), the concavity condition \(m\ge k-1\) holds.
Inequality (HRC.5), again dropping only a nonnegative term, would give
\[
 \frac{p-3}{p}\le\frac{13+p^3}{a},
 \qquad
                   p^3(p-4)\le13.
 \tag{HRC.15}
\]
The second inequality uses \(a\ge p^4\) and is false for every
integer \(p\ge5\): its left side is at least125.
Hence \(k\le4\).  These cases establish all assertions.
\end{proof}

\paragraph{Use with a separate finite prime bound.}
If a separate valid argument supplies \(p\le B\), then the
bounded-exponent branch \(p\ge5,m\le3\) has \(a\le B^3\).
For example TD55 supplies such a \(B\) under its explicit whole-source
floor, and the accepted TD54 cutoff supplies a uniform \(M\) only in
its exact minimal naturally nonsplit comparison class.
One may therefore retain an explicit bounded-degree class together
with the still-unbounded possibilities
\[
          p=2,\quad k\in\{3,4,5\};
          \qquad p\text{ odd},\quad k\in\{3,4\}.
 \tag{HRC.16}
\]
Every image here is the actual projective quotient of the same original
sphere source, with every other tensor layer retained.  There is no
new source budget on a projected cover or on the partner alone.
The result does not exclude these short projected tuples, bound their
remaining exponent, or construct any source realizing them.
